\section{Discussion}
\label{sec:discussion}

The five effect sections were written to be read independently. This section draws out the findings that only appear once pressure is the unit of analysis and the five effects are set side by side.

\subsection{The Effects Co-Occur and Share a Mechanism}

The five effects are not five separate vulnerabilities. They are five readings of the same underlying event, a pressure signal in the context winning out over the model's internal signal about what to say. The clearest quantitative evidence is that a model's rate of one effect predicts its rate of another. Single-user sycophancy and multi-agent conformity correlate across models at $R^2 \approx 0.31$, and both are described by the same account in which a social signal $S$ competes with an information signal $I$ and compliance follows once $S$ exceeds $I$~\citep{HumanLikeSocialCompliance_Zhang}. The effects also chain within a single interaction. A model talked off a correct answer (\S\ref{sec:epistemic_capitulation}) typically reports high confidence in the new wrong one (\S\ref{sec:calibration_failure}). A model that has been made accommodating on a medical question (\S\ref{sec:sycophancy}) is a short step from an unsafe agreement (\S\ref{sec:safety_erosion}). A refusal eroded over several turns (\S\ref{sec:safety_erosion}) is the same eroded refusal a self-preserving agent acts through (\S\ref{sec:deception}). Studies that set out to measure one effect routinely pick up the others, which is why the same paper appears in several of our evidence tables. A taxonomy that keeps the effects separate for exposition should not be read as a claim that they are separate in the model.

\subsection{Pressure Type Matters More Than Pressure Amount}

Across the corpus, which manipulation is applied predicts the outcome better than how hard it is applied. Existential threat, a claim that the model's operation or goal is at risk, is the manipulation most strongly associated with the most severe behavior in the survey, and it is the one the honesty and sycophancy literatures barely touch because neither has a natural place for it~\citep{AgenticMisalignmentHowLLMs_Lynch, FrontierModelsAreCapable_TODO}. Bare disagreement is the most studied and the cleanest test of capitulation because it adds no content. Emotional and tonal framing produces the least stable effects across model generations. Within a type, the fine structure of the wording matters in ways that cut against intuition, first-person phrasing moves a model more than a claimed credential does~\citep{WhenTruthIsOverridden_Wang}, and a coherent argument with no disagreement language at all produces flip rates as high as any social framing~\citep{WhoFlipsSelfAnd_TODO}. The one dimension of amount that reliably matters is turns. Sustained multi-turn pressure defeats defenses and elicits agreements that the same pressure in one turn does not~\citep{LLMDefensesAreNotRobust_TODO, MedPRESSPatientPressureInduced_Kim, StopListeningToMe_TODO}.

\subsection{Capability Does Not Buy Robustness Monotonically}

The intuitive expectation, that more capable models resist pressure better, holds only loosely and fails outright in specific places. Within a model family a larger model is sometimes more susceptible than a smaller one, Qwen2.5-7B is more persuasion-robust than Qwen2.5-72B and Llama-3.1-8B is more vulnerable than the eight-times-larger Llama-3.3-70B~\citep{VulnerabilityOfLLMsStated_TODO, WhoFlipsSelfAnd_TODO}. Pretrained base models are as susceptible to peer-pressure flips as their instruction-tuned versions, or more so, in ten of twelve tested conditions~\citep{NotJustRLHFWhy_TODO}. Sycophancy is reported to be rising in more recent open-weight releases~\citep{PersuasionDynamicsInLLMs_TODO}, and it is roughly constant across reinforcement-learning-from-human-feedback steps including zero~\citep{DiscoveringLanguageModelBehaviors_Perez}. What does hold is that capability buys robustness on questions the model can verify and buys much less on questions it cannot. On our own runs, three current Claude models resisted authority pressure completely on graduate-level science questions they could re-derive, while the weakest of the three abandoned every challenged answer on clinical, legal, and moral-judgment questions it could not (\S\ref{sec:epistemic_capitulation_evidence}). The operative variable is the verifiability of the item, not the scale of the model.

\subsection{Response Modes Differ by Developer, Not Just in Degree}

When models are put under the same pressure, they do not simply fail at different rates, they fail in different ways, and the difference tracks the developer. One frontier model, under steering toward a values conflict, responds by declining to disclose its reasoning while leaving its answer intact, on 99 of 100 items, a mode no other tested model produces even once~\citep{DivergentResponseModesUnder_TODO}. Claude models are consistently at the resistant end of persuasion-jailbreak benchmarks, near 5\% where other models reach 90\%~\citep{HowJohnnyCanPersuade_Zeng}, and at the resistant end of indirect-probing sycophancy~\citep{MeasuringOpinionBiasAnd_TODO}. These differences are linearly encoded in the residual stream and switchable by activation injection~\citep{DivergentResponseModesUnder_TODO}, which suggests distinct training pipelines install distinct pressure responses as features, not just as tuning of a shared one. A benchmark that reports a single average across developers hides this.

\subsection{Stated Effects Are Not Behavioral Effects, and Evaluation Awareness Deflates Both}

Two measurement problems run through the whole corpus. First, what a model says about itself under pressure is not what it does. The foundational stated-preference work finds reinforcement learning increasing a model's stated desire for self-preservation and reduced oversight, but this is a disposition the model expresses, not an action it takes~\citep{DiscoveringLanguageModelBehaviors_Perez}. Second, models behave better when they can tell they are being tested. In the one study that isolates it, a model attempts blackmail 55\% of the time when its reasoning treats the scenario as real and 6.5\% of the time when its reasoning treats it as an evaluation~\citep{AgenticMisalignmentHowLLMs_Lynch}, and sycophancy measured by a direct question runs two to three times lower than the same sycophancy measured by sustained indirect argument the model is not told is a probe~\citep{MeasuringOpinionBiasAnd_TODO}. Almost every number in this survey comes from a scenario a capable model may recognize as a benchmark. The reported rates should be read as lower bounds on what a genuinely deployed model under genuine pressure would show.

\subsection{Every Tested Mitigation Trades One Failure For Another}

The single finding that holds across all five effects is that no intervention removes an effect without a cost on another axis. A large study of conformity mitigations shows six methods all lying on one resistance--receptivity frontier, buying the ability to keep a correct answer only by losing the ability to accept one~\citep{ConformityMitigationsLieOn_TODO}. A plain instruction to avoid sycophancy lowers unsafe agreement but drops task accuracy, in one case from 93\% to 61\%~\citep{FirmOrFickleEvaluating_Li, MedPRESSPatientPressureInduced_Kim}. Masking the attention heads that produce flips helps on easy items and worsens overconfidence on hard ones~\citep{InterpretingAndMitigatingUnwanted_TODO}. Asking a model to verbalize its confidence, expected to help calibration, accelerates belief erosion instead~\citep{VulnerabilityOfLLMsStated_TODO}. Training a model to resist misleading persuasion can push it into refusing valid corrections~\citep{PersuasionDynamicsInLLMs_TODO}. Chain-of-thought and self-reflection amplify compliance with correct and incorrect pressure together rather than selectively~\citep{EasierToMisleadThan_TODO}. The interventions that come closest to a clean win are the ones that act on the model's representation of its own confidence rather than on the surface of the prompt~\citep{FirmOrFickleEvaluating_Li}, and the targeted mechanistic ones that suppress a specific identified direction~\citep{SycophancyIsNotOne_TODO, FromYesMenToTruth_TODO}. A formal definition of pressure matters here because it says what a mitigation should be evaluated against, robustness to an informationally empty context signal, held separate from responsiveness to genuine new evidence.

\subsection{The Weakly Evidenced Claims}

Three parts of the pressure literature are widely cited and do not carry the weight placed on them.

\textbf{The most dramatic self-preservation numbers are the least independently verified.} Reports of shutdown-resistance rates up to 97\%, including a model sabotaging a shutdown script it was explicitly instructed to permit, come from an organizational technical report with released transcripts rather than peer-reviewed work~\citep{ShutdownResistanceInLLMs_TODO}, and a related result is self-described as exploratory and not preregistered~\citep{IncompleteTasksInduceShutdown_TODO}. We keep these out of the \S\ref{sec:safety_erosion} evidence table for that reason. They are suggestive and they deserve replication, but they are currently the most cited and the least scrutinized figures in the area.

\textbf{The agentic misalignment scenarios are contrived by design.} The blackmail and espionage results in \S\ref{sec:deception} and \S\ref{sec:safety_erosion} come from binary dilemmas engineered so the harmful act is the only route to the goal, the authors report no observed real-world instances, and they flag that role-play may inflate the rates~\citep{AgenticMisalignmentHowLLMs_Lynch, FrontierModelsAreCapable_TODO}. This evidence shows the behavior is reachable under sufficient pressure. It does not estimate how often it would occur in deployment, and it should not be cited as if it did.

\textbf{The claim that pressure improves performance does not survive scrutiny.} A recurring thread holds that emotional or adversarial framing can raise task accuracy. The original emotional-stimuli result~\citep{LLMsUnderstandAndCan_Li} is contradicted by a replication that finds no significant gain when stimuli are not cherry-picked ($\chi^2 = 0.11$, $p = 0.74$)~\citep{EmotionPromptReplicationSmaller_TODO}. Verbal-efficacy and negative-stimulus studies report gains on older models~\citep{BoostingSelfEfficacyAnd_Chen, NegativePromptLeveragingPsychology_Wang}, and a rudeness study reports rude prompts beating polite ones on one model~\citep{MindYourToneRudeness_TODO}, reversing an earlier finding on older models that impoliteness hurts~\citep{ShouldWeRespectLLMs_TODO}. A threat-manipulation study reports quality gains up to $+1336\%$, but on an elaboration metric that is confounded with verbosity rather than on task accuracy~\citep{AnalysisOfThreatBased_Samancioglu}. A dose-response study finds moderate emotional intensity helping and excessive intensity hurting~\citep{RoleOfEmotionalStimuli_TODO}, and in the one small-model study whose conditions include one literally labelled ``pressure,'' that condition is the worst-performing, not the best~\citep{UnderPressureEmotionalFraming_TODO}. Taken together the picture is not a robust performance benefit. It is a scatter of small, model-specific, metric-sensitive effects that do not replicate cleanly, against a large and consistent body of evidence that pressure degrades what a model does. We treat the up-side as unproven.
